\documentclass[11pt]{article}
\usepackage{amsmath,amssymb,amsthm,booktabs}
\usepackage[hidelinks]{hyperref}
\emergencystretch=1.5em

\newcommand{\ZZ}{\mathbb Z}
\theoremstyle{plain}
\newtheorem*{theorem}{Theorem}
\newtheorem*{proposition}{Proposition}

\title{Remarks on Bedert's ``On unique sums in Abelian groups''}
\author{Mark Watson}
\date{29 July 2026}

\begin{document}
\maketitle

This note records five small observations about Bedert's paper [B]. The first
concerns the scope of one step in the printed proof for general Abelian groups.
The next two tighten remarks from an earlier version of this note. The last two
record an empty-set convention and finite computations. The theorem discussed
in Section 1 is not disputed.

For a finite subset $A$ of an Abelian group, [B] calls $a_1+a_2$ a unique sum
when its only ordered representations are $(a_1,a_2)$ and $(a_2,a_1)$.

\section{A 2-torsion scope issue}

At [B, line 291], the proof says that a set with no unique sum is balanced:
the sum $a+a$ has another representation, which is taken to have distinct
summands. In a group with 2-torsion, however, the other representation may be
another diagonal pair. Consider
\[
 A=\{0,1,2,3,5,6,8\}\subset \ZZ/12\ZZ.
\]
Here is the complete table of representations, with the order of the two
summands suppressed:
\[
\begin{array}{c|l}
s&\{a,b\}\subset A,\ a+b=s\\ \hline
0&(0,0),(6,6)\\
1&(0,1),(5,8)\\
2&(0,2),(1,1),(6,8)\\
3&(0,3),(1,2)\\
4&(1,3),(2,2),(8,8)\\
5&(0,5),(2,3)\\
6&(0,6),(1,5),(3,3)\\
7&(1,6),(2,5)\\
8&(0,8),(2,6),(3,5)\\
9&(1,8),(3,6)\\
10&(2,8),(5,5)\\
11&(3,8),(5,6)
\end{array}
\]
Every sum has at least two unordered representations, so $A$ has no unique
sum under the definition in [B, line 62]. It is not balanced: neither $0$ nor
$6$ is the midpoint of two distinct members of $A$. In particular, the two
ordered representations of zero are $(0,0)$ and $(6,6)$, so the ordered-count
criterion at [B, lines 118--122] also diverges from the line-62 definition.

The proof of the general-group lower-bound theorem invokes the implication
``no unique sum implies balanced'' at [B, line 458]. Thus the printed proof
needs an odd-order clause at that step, or a separate argument for groups with
2-torsion. This is a scope issue in the proof. It is not a counterexample to
the theorem itself.

The boundary is sharp among cyclic groups of even order.

\begin{theorem}
For $m\geq 0$, there is a set $A\subset \ZZ/(2m)\ZZ$ with no unique sum and
some ordered representation count equal to $2$ if and only if $m\geq 6$.
\end{theorem}

For $m\geq6$, one may take
\[
 \{0,\ldots,m\}\setminus\{m-2\}\ \cup\ \{m+2\}.
\]
The positive family, including the concrete $C_{12}$ witness, and the
exhaustive cases $2m<12$ are machine-checked in \texttt{CycEven.lean}.
\texttt{C12Witness.lean} records the displayed set as a standalone object,
and the full table is recomputed by \texttt{verify.py}.

\section{The {\L}uczak--Schoen upper bound}

{\L}uczak and Schoen [LS, Theorem 2] use the same ordered representation
count as [B]. Taking their parameter $Q=2$ gives, for $p\geq 2^{22}$, a set
$A\subset\ZZ/p\ZZ$ with
\[
 |A|\leq 4(\log_2p)^2,\qquad r_A(g)\geq4\quad(g\in A+A).
\]
Hence $A$ has no unique sum and $m(p)\leq4(\log_2p)^2$. Thus the order of
growth in [B, Theorem 5] already follows from [LS]. Bedert's balanced-set
construction improves the constant and gives a more direct construction.
The author has confirmed this observation and intends to add a note.

\section{A redundant factor in Proposition 1}

For an indexed family $Z=(z_1,\ldots,z_n)$, let $\Sigma(Z)$ be its set of
subset sums and let $d$ be the largest size of a dissociated index set. Then
\[
 |\Sigma(Z)|
 \leq \#\{I\subseteq[n]:I\hbox{ is dissociated}\}
 \leq \sum_{j=0}^{d}\binom nj
 \leq \binom{n+d}{d}.
\]
The proof is the one in the first version of this note. Choose the
colexicographically least index set representing each subset sum. Every set
shattered by this family is dissociated, and Pajor's strengthening of the
Sauer--Shelah lemma completes the count. This removes the first factor from
the bound $\binom nd\binom{n+d}{d}$ in [B, Proposition 1].
\texttt{SubsetSums.lean} checks the first two inequalities, and
\texttt{verify.py} checks the elementary Vandermonde step. The author has
confirmed the improvement and explained that he chose the simpler sufficient
bound because the extra factor does not affect the asymptotic main theorem.

\section{The empty set}

The definition at [B, line 282] makes the empty set balanced, since its
condition is universal over its elements. Proposition 3 at [B, lines
296--298], however, asks for an element $g\in-B$ for every balanced set $B$.
For $B=\varnothing$ there is no such $g$. The proof at [B, lines 349--350]
also uses that a balanced set contains at least two distinct elements.

The harmless repair is to require balanced sets to be nonempty in the
definition, or to state the propositions for nonempty balanced sets. With
that convention, the argument at [B, lines 349--350] supplies the stronger
fact that such a set has at least two elements.

\section{Exact values of $m(p)$}

The current computations give:
\[
\begin{array}{c|rrrrrrrrrrrrrrrr}
p&3&5&7&11&13&17&19&23&29&31&37&41&43&47&53&59\\ \hline
m(p)&3&4&5&7&7&8&9&10&11&11&12&13&13&13&14&15
\end{array}
\]
The verification tiers are not uniform:
\begin{center}
\begin{tabular}{@{}p{0.33\textwidth}|p{0.61\textwidth}@{}}
$p$&verification tier\\ \hline
$3,5,7,11,13,17,19,$\newline $23,29,31,37$&
exact solver output, independently corroborated; not certified here\\
$41,43,47$&certified optimality and verified witness\\
$53,59$&exact by exhaustion and verified witness; not certified
\end{tabular}
\end{center}
Here ``certified'' means that a retained machine-checkable certificate was
verified. The values at $53$ and $59$ have exhaustive lower-bound searches
but no DRAT certificate. The checking script redoes the small cases feasible
in pure Python and verifies the displayed witnesses at $41,43,47,53,59$.

\section*{Checking}

Run
\[
\texttt{lake exe cache get \&\& lake build}
\qquad\text{and}\qquad
\texttt{python3 verify.py}.
\]
The printed axiom sets for the main Lean theorems should contain exactly
\texttt{propext}, \texttt{Classical.choice}, and \texttt{Quot.sound}.

\begin{thebibliography}{LS}
\bibitem[B]{bedert} B.~Bedert, \emph{On unique sums in Abelian groups},
Combinatorica \textbf{44}(2) (2024), 269--298; arXiv:2303.15134.
\bibitem[LS]{ls} T.~{\L}uczak and T.~Schoen, \emph{On a problem of Konyagin},
Acta Arith.\ \textbf{134}(2) (2008), 101--109.
\bibitem[P]{pajor} A.~Pajor, \emph{Sous-espaces $l_1^n$ des espaces de Banach},
Travaux en Cours 16, Hermann, Paris, 1985.
\end{thebibliography}

\end{document}
